% Options for packages loaded elsewhere
% Options for packages loaded elsewhere
\PassOptionsToPackage{unicode}{hyperref}
\PassOptionsToPackage{hyphens}{url}
\PassOptionsToPackage{dvipsnames,svgnames,x11names}{xcolor}
%
\documentclass[
  letterpaper,
  DIV=11,
  numbers=noendperiod]{scrartcl}
\usepackage{xcolor}
\usepackage{amsmath,amssymb}
\setcounter{secnumdepth}{-\maxdimen} % remove section numbering
\usepackage{iftex}
\ifPDFTeX
  \usepackage[T1]{fontenc}
  \usepackage[utf8]{inputenc}
  \usepackage{textcomp} % provide euro and other symbols
\else % if luatex or xetex
  \usepackage{unicode-math} % this also loads fontspec
  \defaultfontfeatures{Scale=MatchLowercase}
  \defaultfontfeatures[\rmfamily]{Ligatures=TeX,Scale=1}
\fi
\usepackage{lmodern}
\ifPDFTeX\else
  % xetex/luatex font selection
\fi
% Use upquote if available, for straight quotes in verbatim environments
\IfFileExists{upquote.sty}{\usepackage{upquote}}{}
\IfFileExists{microtype.sty}{% use microtype if available
  \usepackage[]{microtype}
  \UseMicrotypeSet[protrusion]{basicmath} % disable protrusion for tt fonts
}{}
\makeatletter
\@ifundefined{KOMAClassName}{% if non-KOMA class
  \IfFileExists{parskip.sty}{%
    \usepackage{parskip}
  }{% else
    \setlength{\parindent}{0pt}
    \setlength{\parskip}{6pt plus 2pt minus 1pt}}
}{% if KOMA class
  \KOMAoptions{parskip=half}}
\makeatother
% Make \paragraph and \subparagraph free-standing
\makeatletter
\ifx\paragraph\undefined\else
  \let\oldparagraph\paragraph
  \renewcommand{\paragraph}{
    \@ifstar
      \xxxParagraphStar
      \xxxParagraphNoStar
  }
  \newcommand{\xxxParagraphStar}[1]{\oldparagraph*{#1}\mbox{}}
  \newcommand{\xxxParagraphNoStar}[1]{\oldparagraph{#1}\mbox{}}
\fi
\ifx\subparagraph\undefined\else
  \let\oldsubparagraph\subparagraph
  \renewcommand{\subparagraph}{
    \@ifstar
      \xxxSubParagraphStar
      \xxxSubParagraphNoStar
  }
  \newcommand{\xxxSubParagraphStar}[1]{\oldsubparagraph*{#1}\mbox{}}
  \newcommand{\xxxSubParagraphNoStar}[1]{\oldsubparagraph{#1}\mbox{}}
\fi
\makeatother


\usepackage{longtable,booktabs,array}
\newcounter{none} % for unnumbered tables
\usepackage{calc} % for calculating minipage widths
% Correct order of tables after \paragraph or \subparagraph
\usepackage{etoolbox}
\makeatletter
\patchcmd\longtable{\par}{\if@noskipsec\mbox{}\fi\par}{}{}
\makeatother
% Allow footnotes in longtable head/foot
\IfFileExists{footnotehyper.sty}{\usepackage{footnotehyper}}{\usepackage{footnote}}
\makesavenoteenv{longtable}
\usepackage{graphicx}
\makeatletter
\newsavebox\pandoc@box
\newcommand*\pandocbounded[1]{% scales image to fit in text height/width
  \sbox\pandoc@box{#1}%
  \Gscale@div\@tempa{\textheight}{\dimexpr\ht\pandoc@box+\dp\pandoc@box\relax}%
  \Gscale@div\@tempb{\linewidth}{\wd\pandoc@box}%
  \ifdim\@tempb\p@<\@tempa\p@\let\@tempa\@tempb\fi% select the smaller of both
  \ifdim\@tempa\p@<\p@\scalebox{\@tempa}{\usebox\pandoc@box}%
  \else\usebox{\pandoc@box}%
  \fi%
}
% Set default figure placement to htbp
\def\fps@figure{htbp}
\makeatother





\setlength{\emergencystretch}{3em} % prevent overfull lines

\providecommand{\tightlist}{%
  \setlength{\itemsep}{0pt}\setlength{\parskip}{0pt}}



 


\KOMAoption{captions}{tableheading}
\makeatletter
\@ifpackageloaded{caption}{}{\usepackage{caption}}
\AtBeginDocument{%
\ifdefined\contentsname
  \renewcommand*\contentsname{Table of contents}
\else
  \newcommand\contentsname{Table of contents}
\fi
\ifdefined\listfigurename
  \renewcommand*\listfigurename{List of Figures}
\else
  \newcommand\listfigurename{List of Figures}
\fi
\ifdefined\listtablename
  \renewcommand*\listtablename{List of Tables}
\else
  \newcommand\listtablename{List of Tables}
\fi
\ifdefined\figurename
  \renewcommand*\figurename{Figure}
\else
  \newcommand\figurename{Figure}
\fi
\ifdefined\tablename
  \renewcommand*\tablename{Table}
\else
  \newcommand\tablename{Table}
\fi
}
\@ifpackageloaded{float}{}{\usepackage{float}}
\floatstyle{ruled}
\@ifundefined{c@chapter}{\newfloat{codelisting}{h}{lop}}{\newfloat{codelisting}{h}{lop}[chapter]}
\floatname{codelisting}{Listing}
\newcommand*\listoflistings{\listof{codelisting}{List of Listings}}
\makeatother
\makeatletter
\makeatother
\makeatletter
\@ifpackageloaded{caption}{}{\usepackage{caption}}
\@ifpackageloaded{subcaption}{}{\usepackage{subcaption}}
\makeatother
\usepackage{bookmark}
\IfFileExists{xurl.sty}{\usepackage{xurl}}{} % add URL line breaks if available
\urlstyle{same}
\hypersetup{
  pdftitle={Time translations},
  colorlinks=true,
  linkcolor={blue},
  filecolor={Maroon},
  citecolor={Blue},
  urlcolor={Blue},
  pdfcreator={LaTeX via pandoc}}


\title{Time translations}
\author{}
\date{}
\begin{document}
\maketitle


In class we learned that the (linear) momentum \(p\) is the
``generator'' of translations in space. (We work in one spacial
dimension for simplicity.) Concretely what this meant is that we can
construct from \(p\) an operator \(U\) on the Hilbert space of functions
\(\{f(x)\}\) of the coordinate \(x\) so that \[
U(a)f(x) = f(x+a)
\] Recall that since we want to preserve the Hilbert space norm, \(U\)
should be unitary. From this and certain limits, we worked out that
\(U(a)\) could be expressed in terms of the absolutely convergent sum
for the exponential \[
\exp \left( \frac i\hbar ap \right) = \sum_{n=0}^\infty \frac1{n!} \left( \frac i\hbar ap \right)^n
\] In other words \(U(a)=e^{a \frac {d}{dx} } f(x) = f(x+a)\).

\begin{quote}
\begin{enumerate}
\def\labelenumi{\arabic{enumi}.}
\setcounter{enumi}{-1}
\tightlist
\item
  Show this explicity by expanding out the exponential and
  reinterpreting the result as a Taylor series.
\end{enumerate}
\end{quote}

The exponential sum is \[
e^{ a\frac d{dx} } = \sum_{n=0}^\infty \frac 1{n!} a^n \frac {d^n}{dx^n}
\] Let \(f(x)\) be any smooth function and Taylor expand \(f(x+a)\)
around some constant \(a\): \[
f(x+a)=\sum_{n=0}^\infty \frac 1{n!} a^n f^{(n)}
\] where \(f^{(n)}\) is a common notation for the \(n\)th derivative of
\(f\). But then we are done, because these two expansions are identical.

In this set of notes, we will work out the same set of steps for
translations in time. That is, we aim to find the generator of time
translations. Time in Newtonian physics is nothing like space. All
observers agree on the passing of time, and at any particular moment in
time, space is Euclidean. Coordinates like \(x\) are merely choices we
make to parameterize this space. Time on the other hand defines
\emph{dynamics} though Newton's Second Law. Because of this, what we
will find that the rabbit hole goes far deeper in time\ldots{}

So, by complete analogy, we are looking for an operator \(U(c)\) that
maps functions of time \(f(t)\mapsto f(t+c)\) for any constant \(c\). By
all the same arguments, this has to have the form
\(U(c)=e^{c \frac d{dt}}\) which we can write as
\(e^{-\frac i\hbar ch }\) where \(h\) is a Hermitian operation on
functions called the generator of time translations. Just as \(p\) was
the generator of translations, we now need to identify this operator.

\begin{quote}
\begin{enumerate}
\def\labelenumi{\arabic{enumi}.}
\tightlist
\item
  Find the dimensions of the generator \(h\). What real (\emph{i.e.} has
  no imaginary part) quantity do you know with these dimensions?
\end{enumerate}
\end{quote}

Let \([A]\) denote ``the dimensions of \(A\)''. Exponentiation of a
dimensionful quantity is meaningless (ask, if you don't understand
this!), so \([\frac i\hbar ch]=1\). The imaginary unit is dimensionless
since it is just a number.
\([\hbar]=[x]\cdot [p]= L \cdot ML/T=ML^2/T\). Since \(c\) enters in
\(f(t)\mapsto f(t+c)\), \([c]=[t]=T\). Thus
\(1=[\frac {ch}\hbar ]=T/(ML^2) \cdot T \cdot [h]\), which implies
\([h]=ML^2/T^2\). But these are the dimensions of mass times
velocity-squared, which is an energy. So whatever \(h\) is, it must have
the dimesions of energy.

Now to identify what \(h\) actually is, we use the observation that the
generator of time translations is what we use to write Newton's second
law \[
F=\frac {dp}{dt}
\] and it is unique in that way. We must also recall that the momentum
is \emph{defined} as a derivative of the kinetic energy \(K\).

\begin{quote}
\begin{enumerate}
\def\labelenumi{\arabic{enumi}.}
\setcounter{enumi}{1}
\tightlist
\item
  If you did not derive this in class, derive explicitly from the
  example of a free classical particle of mass \(m\) and velocity
  \(\bf v\) a formula for \(\bf p\) as a derivative of \(K\).
\end{enumerate}
\end{quote}

The energy of a free particle is \(K=\frac12 m v^2\) where
\(v=\sqrt{{\bf v} \cdot {\bf v}}\). Thus, component-by-component writing
\({\bf v}=(v^i)\) for the 3 differnt directions \(i=x,y,z\),
\(\frac{\partial K}{\partial v^i}=2\times \frac12 m v_i= mv_i=p_i\).
Thus, component-by-component \[
{\bf p} = \frac{\partial K}{\partial {\bf v}}
\] We can take this formula as the definition of \({\bf p}\). We say
that ``\({\bf p}\) is the momentum canonically conjugate to
\({\bf x}\)''. (More on this terminology below.)

Besides the kinetic energy there is another potential energy for any
conservative force: If \({\bf R}^3\),
\(0=\int_\gamma {\bf F}\cdot d{\bf x}=0\) for any closed contour
\(\gamma\), then there is a function \(U({\bf x})\) such that
\({\bf F}=-{\nabla}U\). Importantly, although \(K=K(x,v)\) is allowed to
depend on \(x\), \(U(x)\) is \emph{not} allowed to depend on \(v\). (I
know, unfair right?) Now consider the total energy function
\(H=K(x,v)+U(x)\).

\begin{quote}
\begin{enumerate}
\def\labelenumi{\arabic{enumi}.}
\setcounter{enumi}{2}
\tightlist
\item
  For this problem, treat \(p\) and \(v\) as independent variables.\\
\end{enumerate}

\begin{enumerate}
\def\labelenumi{\alph{enumi}.}
\tightlist
\item
  Show that for any kinetic energy, the combination
  \[H(x,p)=pv-K(x,v)+U(x)\] is independent of \(v\).\\
\item
  Assume that \(K(x,v)=f(x) v^\alpha\) for some real number \(\alpha\).
  Show that \(pv-K=(\alpha -1)K\).\\
\item
  Combine parts a. and b. to conclude that, if \(K\) is quadratic in
  \(v\), then \(H(x,p)=K(x,v)+U(x)\).
\end{enumerate}
\end{quote}

3 a. Just calculate the partial derivative and use the definition of the
momentum: \[
\frac{\partial H}{\partial v}
=\frac{\partial }{\partial v}\left[ pv-K(x,v)+U(x) \right]
=p-\frac{\partial H}{\partial v}
=p-p
=0
\]

\begin{enumerate}
\def\labelenumi{\alph{enumi}.}
\setcounter{enumi}{1}
\item
  If \(K=f(x)v^\alpha\), then
  \(\frac{\partial K}{\partial v}=\alpha f(x)v^{\alpha-1}\) which
  implies that \(v\frac{\partial K}{\partial v}=\alpha K\). But this is
  the same as \(pv=\alpha K\), so \(pv-K=(\alpha -1)K\).
\item
  If \(K\) is quadratic in \(v\), then \(\alpha=2\). When that happens,
  \(pv-K=(\alpha -1)K=K\), so \(H(x,p)=pv-K(x,v)+U(x)=K(x,v)+U(x)\).
\end{enumerate}

This took a lot of effort that seems like it should not have been
necessary, so we should go back an look at how \(K\) was originally
defined.

Recall that all these energy consideratoins came from work without
dissipative forces. Then \(F=-\nabla U\) for some \(U\) and
\(\int_\gamma F\cdot dx = - \int_\gamma \nabla U\cdot dx=\left. U\right|_{-\partial \gamma}\).
But
\(\int_\gamma F\cdot dx = \int_\gamma \dot p \cdot dx = \int_\gamma m \dot v \cdot v dt=\frac m2 \int_\gamma\frac {dv^2}{dt} dt=\left.\frac m2 v^2\right|_{\partial \gamma}\).
Thus \(\left[\frac m2 v^2+U\right]_{\partial \gamma}=0\) and this holds
along any path \(\gamma\) in space. Therefore, \(\frac m2 v^2+U =\)
constant. This was the original motivation for introducing
\(K=\frac m2 v^2\).

Now in general \begin{align}
\int_\gamma F\cdot dx &= \int_\gamma \dot p \cdot dx 
= \int_\gamma \frac d{dt}\left(\frac{\partial K}{\partial v}\right) \cdot dx 
= \int_\gamma \frac d{dt}\left(\frac{\partial K}{\partial v}\right) \cdot vdt \cr
&= \int_\gamma \frac d{dt}\left(\alpha f(x)v^{\alpha -1}\right) \cdot vdt
= \int_\gamma \frac d{dt}\left(f(x)v^\alpha\right)dt
= \int_\gamma \frac {dK}{dt}dt
=\left.K\right|_\gamma
\end{align} so we, again recover the work-kinetic energy theorem, but
this time for any kinetic energy of the form \(K=f(x)v^\alpha\).

\subsection{Hamiltonian formulation of
dynamics}\label{hamiltonian-formulation-of-dynamics}

But wait, what?! \(H\) clearly depends on \(v\)! What's going on??

\begin{quote}
\begin{enumerate}
\def\labelenumi{\arabic{enumi}.}
\setcounter{enumi}{3}
\tightlist
\item
  Rewrite the kinetic energy part of the Hamiltonian \(H\) for a
  particle of mass \(m\) in terms of momentum. Use this expression to
  show that \[
  \frac{d x}{dt} = \frac{\partial H}{\partial p}
  \] Speculate about what this means for the dependence of the total
  energy on position, velocity, and linear momentum.
\end{enumerate}
\end{quote}

The kinetic energy of a particle of mass \(m\) is \(K=\frac m2 v^2\).
Its momentum is is \(p=\partial K/\partial v=mv\), to \(v=\frac pm\),
and \(K=\frac12 pv=\frac1{2m}p^2\). Thus
\(dx/dt=v=\partial K/\partial p\), but since the potential energy does
not depend on \(v\), it cannot depend on \(p\). Therefore,
\(dx/dt=\partial H/\partial p\).

\begin{quote}
\begin{enumerate}
\def\labelenumi{\arabic{enumi}.}
\setcounter{enumi}{5}
\tightlist
\item
  For a partcile of mass \(m\) in Euclidean space, show that Newton II
  can be re-expressed as \[
  \frac{d p}{dt} = -\frac {\partial H}{\partial x}
  \]
\end{enumerate}
\end{quote}

This is again just a calculation: \[
\frac{\partial H}{\partial x}=\frac{\partial K}{\partial x}+\frac{\partial U}{\partial x}=-F=-\frac{d p}{dt}
\]

Congratulations! You have derived Hamiltons ``canonical equations''.
Let's recap. The total energy function \(H(x,p)\) is a function on the
\textbf{phase space} of positions and momenta. This space is
\(2n\)-dimensional when the \textbf{configuration space} (a.k.a.
``space'', \emph{e.g.} \({\bf R}^3\)) is \(n\)-dimensional. In this
formulation of Newtonian physics, particle dynamics corresponds to the
allowed paths in phase space. If there is no \emph{explicit}
time-dependence in the system (\emph{i.e.} the system is not being
driven with time-dependent forces), then all intersting dynamical
quantities can be formulated as functions on this phase space. So all
the quantities we care about are phase space functions. (We will return
to explicit time-dependence later.)

In this phase space formulation, the dynamics of any quantity \(f(x,p)\)
is encoded in Hamilton's equations. That is, the ``Hamiltonian
(function)'' completely determines the dynamics of the system.

\begin{quote}
\begin{enumerate}
\def\labelenumi{\arabic{enumi}.}
\setcounter{enumi}{6}
\tightlist
\item
  Let \(f(x,p)\) be any differentiable function on the phase space. Use
  the chain rule and Hamilton's equations to write a formula for
  \(df/dt\).
\end{enumerate}
\end{quote}

\begin{align}
\frac{df}{dt} &= \frac{\partial f}{\partial x}\frac{dx}{dt} + \frac{\partial f}{\partial p}\frac{dp}{dt} 
=\frac{\partial f}{\partial x}\frac{\partial H}{\partial p} - \frac{\partial f}{\partial p}\frac{\partial H}{\partial x}
\end{align}

\subsubsection{Poisson Bracket}\label{poisson-bracket}

Looking at this formula we found in problem 6, we notice an interesting
structure that encourages us to introduce the \textbf{Poisson bracket}:
For any two phase space functions \(f\) and \(g\), define the function
\[
\{f, g\}=\frac{\partial f}{\partial p}\frac{\partial g}{\partial x} - \frac{\partial f}{\partial x}\frac{\partial g}{\partial p}
\]

\begin{quote}
\begin{enumerate}
\def\labelenumi{\arabic{enumi}.}
\setcounter{enumi}{7}
\tightlist
\item
  Rewrite your dynamical equation in terms of the Poisson bracket.
\end{enumerate}
\end{quote}

Almost by inspection, the equation can be written as \[
\frac {df}{dt} = \{H,f\}
\]

It is important to note that we did not introduce any new structure into
the theory ``by hand''. (An example of a structure that is just thrown
in by hand is the cross product in particle mechanics.) Even the dot
product is not a native thing to the vector space \({\bf R}^3\), but
added as an assumtion that makes the vector space into an ``inner
product space'' we can Euclidean. The Poisson bracket, by contrast, is a
stucture that comes with any phase space \emph{automatically}. It was
handed to us by the math Gods as a gift. This is the reason
mathematicians call it ``canonical'', and this particular type of
canonical structure is called ``symplectic'' in the math jargon.

\begin{quote}
\begin{enumerate}
\def\labelenumi{\arabic{enumi}.}
\setcounter{enumi}{8}
\tightlist
\item
  For any 3 phase space functions \(f\), \(g\), and \(h\), check that
  the Poisson bracket satisfies the following properties:\\
\end{enumerate}

\begin{enumerate}
\def\labelenumi{\roman{enumi}.}
\tightlist
\item
  \(\{g, f+h\} = \{f, g\}+ \{f, h\}\) and for any real number
  \(c\in {\bf R}\), \(\{f, cg\}=c\{f, g\}\)\\
\item
  \(\{g, f\} = -\{f, g\}\)\\
\item
  \(\{f, \{g, h\}\} + \{g,\{h, f\}\} + \{h, \{f, g\}\} = 0\)
\end{enumerate}
\end{quote}

The first property just follows from how derivatives behave. The second
property is obvious. The third is an annoying calculation: \begin{align}
\{h,\{f,g\}\}&=
\left\{h, \frac{\partial f}{\partial p}\frac{\partial g}{\partial x} - \frac{\partial f}{\partial x}\frac{\partial g}{\partial p} \right\}\cr
&=\left\{h, \frac{\partial f}{\partial p}\right\}\frac{\partial g}{\partial x} - \left\{ h, \frac{\partial f}{\partial x}\right\}\frac{\partial g}{\partial p}
+\frac{\partial f}{\partial p}\left\{h, \frac{\partial g}{\partial x}\right\} - \frac{\partial f}{\partial x}\left\{ h, \frac{\partial g}{\partial p} \right\}
\end{align} Since this is such a mess, we simplify the notation by
writing the derivatives as subscripts: \begin{align}
\{h,\{f,g\}\}&=\left\{h, f_p\right\}g_x- \left\{ h, f_x\right\}g_p +f_p\left\{h, g_x\right\} - f_x\left\{ h, g_p \right\}\cr
&=
h_p f_{xp}g_x - h_xf_{pp}g_x - h_p f_{xx}g_p+h_x f_{px}g_p\cr
&+f_p h_p g_{xx} - f_p h_x g_{px} - f_x h_p g_{xp}+f_x h_xg_{pp}\cr
&=
f_{xx}\left(- g_p h_p \right)
+f_{xp}\left(g_ph_x +  g_x h_p\right)
+f_{pp}\left(- g_x h_x\right)\cr
&+g_{xx}\left(h_p f_p \right)  
+g_{xp}\left( - h_x f_p  - h_pf_x \right)
+g_{pp}\left(h_xf_x \right)
\end{align} Now we need to add the cyclic permutations
\(f\to g\to h\to f\), but here is what will happen: Note that if we do
this permutation to the first line of the result, we get the second line
with all the signs switched. This means that when we add up the
permutations, the second line gets canceled by the first line of the
next permutation, and this repeats, so everything cancels when we add up
all three.

If that was too slick for you, I'll do it explicitly: \begin{align}
&f_{xx}\left(- g_p h_p \right)
+f_{xp}\left(g_ph_x +  g_x h_p\right)
+f_{pp}\left(- g_x h_x\right)\cr
&+g_{xx}\left(h_p f_p \right)  
+g_{xp}\left( - h_x f_p  - h_pf_x \right)
+g_{pp}\left(h_xf_x \right)\cr~\cr
&+g_{xx}\left(- h_p f_p \right)
+g_{xp}\left(h_pf_x +  h_x f_p\right)
+g_{pp}\left(- h_x f_x\right)\cr
&+h_{xx}\left(f_p g_p \right)  
+h_{xp}\left( - f_x g_p  - f_pg_x \right)
+h_{pp}\left(f_xg_x \right)\cr~\cr
&+h_{xx}\left(- f_p g_p \right)
+h_{xp}\left(f_pg_x +  f_x g_p\right)
+h_{pp}\left(- f_x g_x\right)\cr
&+f_{xx}\left(g_p h_p \right)  
+f_{xp}\left( - g_x h_p  - g_p h_x \right)
+f_{pp}\left(g_xh_x \right)
\end{align} Now we finally see that all the terms enter twice with
opposite signs, so they cancel pairwise.

\subsection{Recapitulation}\label{recapitulation}

The fundamental equations of Newtonian dynamics can be reformulated as
statements about operations on a linear space \(\mathscr H\) that almost
magically comes with a Poisson bracket \(\{\cdot, \cdot\}\), that
promotes it to a very special sort of algebra that is intimately tied to
a specific class of differential equations and their solutions. In this
formalism, time translations are generated by a distinguished function
called the Hamiltonian.

This space is not just a vector space with an norm (\emph{i.e.} the
Hilbert space part), but a Lie algebra because it automatically comes
with a Poisson bracket that satisfies the defining relations of a Lie
bracket.

\subsection{Time dependence}\label{time-dependence}

Hamilton's formulation of dynamics shows that the time evolution of any
function \(f\in {\mathscr H}\), that is any function \(f\) on the phase
space of \(p\)s and \(x\)s is generated, through the Poisson Bracket, by
the Hamiltonian. This, however, is not the most general behavior because
it is always possible to have explicit dependence on time. To describe
this situation, we talk about the \emph{extended} phase space of \(x\)s,
\(p\)s, and \(t\).

\begin{quote}
\begin{enumerate}
\def\labelenumi{\arabic{enumi}.}
\setcounter{enumi}{9}
\tightlist
\item
  \begin{enumerate}
  \def\labelenumii{\alph{enumii}.}
  \tightlist
  \item
    Show that for a function on the extended phase space \[
    \frac {df}{dt} = \frac {\partial f}{\partial t} + \{f,H\}
    \]
  \end{enumerate}
\end{enumerate}

\begin{enumerate}
\def\labelenumi{\alph{enumi}.}
\setcounter{enumi}{1}
\tightlist
\item
  Use this to show that, unless the Hamiltonian \(H\) depends explicitly
  on time, it is \textbf{conserved}, where ``conserved'' means constant
  in time.
\end{enumerate}
\end{quote}

9 a. A function on the extended phase space has the form \(f(x, p, t)\).
Therefore, \(\frac {df}{dt} = \frac {\partial f}{\partial t}+\) parts
already there when it does not depend explicitly on \(t\). But we
already know these are generated by \(H\) through the Poisson bracket.
This gives the stated equation.\\
b. Plugging in \(f\to H\) gives \[
\frac {dH}{dt} = \frac {\partial H}{\partial t} + \{H,H\}
= \frac {\partial H}{\partial t}
\] The \(\{H, H\}=-\{H, H\}\) term vanishes by antisymmetry of P.B., so
\(\frac {dH}{dt}=0\) iff \(\frac {\partial H}{\partial t}=0\).

\section{Hamiltonian quantization}\label{hamiltonian-quantization}

\begin{quote}
\begin{enumerate}
\def\labelenumi{\arabic{enumi}.}
\setcounter{enumi}{10}
\tightlist
\item
  Suppose we have an algebra of operators \(\mathscr O\) with a
  multiplication that may or may not be associative: For any \(f\),
  \(g\), \(h\in \mathscr O\) \[
  (h(fg)) \stackrel?= ((hf)g)
  \] and define the \textbf{commutator} \[
  [f,g] = fg - gf
  \] Show that \([\cdot, \cdot]\) satisfies the Jabobi identity if the
  algebra multiplication is associative.
\end{enumerate}
\end{quote}

Let's look at the first of three terms \[
[h,[f,g]]=(h(fg)) - ((fg)h) + ((gf)h) - (h(gh))
\] The other two terms are the permulations \(f\to g\to h\to f\) and
\(f\to h\to g\to f\).

An example of an associative algebra is the algebra of matrices with
matrix multiplication.




\end{document}
